\section{Hessian摂動の局所特異性とPfaffian graph上の解消}
\label{sec:pfaffian-hessian-local-geometry}
本節は2026年11月1日のPfaffian境界原稿（旧論文III）の
第23--24節に含まれる結果を、証明のある現行TeXとして再構成する。
旧稿の普遍非線形接触分類（第21--22節）と
全次数の頂点核の決定（第10節）はまだ本文を照合・統合していない。
従って本節は、実際のBrill--Noether境界が下記の
正規形と接触同値であることを新たに証明するものではない。
固定した滑らかな三次式と係数$c$の模型についての定理である。

\subsection{五変数Hessian摂動のMilnor環とTjurina環}
$V$を五次元複素ベクトル空間とし、
$F\in\operatorname{Sym}^3V^*$が滑らかな
$X=V_+(F)\subset\PP(V)$を定めるとする。
座標$x_0,\ldots,x_4$を選び
\[
 H_F=\det\left(\frac{\partial^2 F}{\partial x_i\partial x_j}\right),
 \qquad g_{F,c}=F+cH_F,\qquad c\in\CC
\]
と置く。$H_F$は五次斉次式である。

\begin{lemma}[正則な初期式によるfiltered Koszul比較]
\label{lem:hessian-filtered-koszul}
$R=\CC[[x_0,\ldots,x_4]]$、$\mathfrak m=(x_0,\ldots,x_4)$とし、
$b_i\in\mathfrak m^2$の二次初期式$b_i^{(2)}$が
$\CC[x_0,\ldots,x_4]$の正則列をなすとする。
$J=(b_0,\ldots,b_4)$なら
\[
 \operatorname{gr}_{\mathfrak m}(R/J)
 \simeq\CC[x_0,\ldots,x_4]/(b_0^{(2)},\ldots,b_4^{(2)}),
 \qquad J\cap\mathfrak m^{k+2}=\mathfrak m^kJ\quad(k\ge0).
\]
\end{lemma}
\begin{proof}
$v=\sum_i a_ib_i$の係数の最小次数を$r$とする。
次数$r+2$の斉次項が消えるなら、
$\sum_i a_i^{(r)}b_i^{(2)}=0$である。
正則列の斉次syzygyはKoszul syzygyで生成される。
このKoszul syzygyを$b_i$に持ち上げた零関係式を引けば、
$v$を保ったまま係数の最小次数を上げられる。
$\mathfrak m$進完備性と逐次の次数上昇を使って、
$v\in\mathfrak m^{k+2}$の表示係数を$\mathfrak m^k$に選べる。
従って交叉公式が得られる。同じ手順で初期イデアルが
$(b_i^{(2)})$と一致するので、associated gradedの同型も従う。
\end{proof}

\begin{proposition}[Hessian摂動のMilnor数とTjurina数]
\label{prop:hessian-milnor-tjurina}
任意の$c\in\CC$について$g_{F,c}$は原点に孤立臨界点を持ち、
Milnor数とTjurina数は
\[
 \mu(g_{F,c})=32,\qquad
 \tau(g_{F,c})=
 \begin{cases}32&c=0,\\31&c\ne0.\end{cases}
\]
Milnor環の$\mathfrak m$進Hilbert級数は$(1+t)^5$、
$c\ne0$のTjurina環のそれは
$1+5t+10t^2+10t^3+5t^4$である。
従って$c\ne0$の五変数超曲面芽は斉次三次錐と
接触同値でない。
\end{proposition}
\begin{proof}
$F_i=\partial_iF$の共通の射影零点があれば、
Eulerの等式から$F$の特異点を与えてしまう。
従って五つの二次式$F_i$は正則列であり、
\[
 M_F=\CC[x_0,\ldots,x_4]/(F_0,\ldots,F_4),\qquad
 \operatorname{Hilb}(M_F;t)=(1+t)^5.
\]
次数五のsocleは一次元で、$[H_F]$がその非零生成元である。
これは$(F_i)$と$(x_i)$のKoszul分解を
$F_i=\frac12\sum_jF_{ij}x_j$によって比較すれば、
最上位Torの生成元が$2^{-5}H_F$へ移ることからも確認できる。

$g_i=\partial_i g_{F,c}=F_i+c\,\partial_iH_F$の
二次初期式は$F_i$なので、
補題\ref{lem:hessian-filtered-koszul}により
\[
 \operatorname{gr}_{\mathfrak m}
 \bigl(R/(g_0,\ldots,g_4)\bigr)\simeq M_F.
\]
したがってMilnor環の次元は32で、
$[H_F]$はなお非零の一次元socleを生成する。
Eulerの公式からMilnor環において
\[
 3[F]+5c[H_F]=0,\qquad
 [g_{F,c}]=-\frac{2c}{3}[H_F].
\]
$c=0$なら$[g_{F,c}]=0$であり、
$c\ne0$なら$[g_{F,c}]$は非零socle元となる。
Tjurina環はこの元でMilnor環を割ったものなので
表示した次元とHilbert級数が従う。
接触同値はTjurina環を保つため、最後の主張も従う。
\end{proof}

\begin{remark}[係数の二つの軌道]
$c,c'\ne0$なら$\lambda^2c=c'$を満たす$\lambda$に対して
$\lambda^{-3}g_{F,c}(\lambda x)=g_{F,c'}(x)$。
従ってこの固定$F$の五変数模型で$c=0$と$c\ne0$は
二つの接触同値類となる。この事実だけから、
十変数のPfaffian境界の全接触分類は導けない。
\end{remark}

\subsection{Pfaffian graphと相対次元五の二次ファイブレーション}
$W=\bigwedge^2V$を十次元空間とし、
五つの$4\times4$ Pfaffianからなる写像を
$q:W\to V^*$（体積形式による同一視を含む）とする。
$\Sigma=\{q=0\}$をrank二以下の錐とし、
\[
 {\cal H}=\operatorname{Bl}_{\Sigma}W
 \simeq \operatorname{Tot}_{\PP^4}(\bigwedge^2\mathcal Q)
\]
をPfaffian graphの標準的な模型とする。
$\mathcal Q$は階数四の普遍商束である。
基底の斉次代表$p$と束内のPfaffian二次式$r$を使うと
$q(a)=rp$であり、$p'=up$に対して$r'=u^{-1}r$となる。

\begin{proposition}[Pfaffian境界模型のstrict transform]
\label{prop:pfaffian-hessian-strict-transform}
$D_{F,c}=V(F(q(a))+cH_F(q(a)))\subset W$とおく。
その$\mathcal H$上のstrict transformは
\[
 \widetilde D_{F,c}
 =\{F(p)+cr^2H_F(p)=0\}\subset\mathcal H.
\]
これは適切な$\PP^4$上の三次線束の大域切断で定義され、
零切断上の$X$の近傍で滑らかである。
\end{proposition}
\begin{proof}
斉次性により
\[
 F(rp)+cH_F(rp)
 =r^3\bigl(F(p)+cr^2H_F(p)\bigr).
\]
例外因子$r=0$が括弧の零集合の全成分とはならないから、
三重の例外因子を除いた表示がstrict transformを与える。
$p'=up,\ r'=u^{-1}r$の変換で括弧は$u^3$倍される。
零切断$r=0$では方程式の微分が$dF$に一致する。
$X$の滑らかさから$dF\ne0$であり、結論が従う。
\end{proof}

\begin{proposition}[四次元孤立特異点の解消とのファイバー積]
\label{prop:pfaffian-hessian-quadric-fibration}
$Y_{F,c}=V(g_{F,c})\subset\CC^5$とし、
$\mathcal L=\operatorname{Bl}_0\CC^5
=\operatorname{Tot}_{\PP^4}\mathcal O(-1)$とする。
線束座標$t$により
\[
 \widetilde Y_{F,c}=\{F(p)+ct^2H_F(p)=0\}
 \subset\mathcal L
\]
と書けば、$X$の近傍で$\widetilde Y_{F,c}$は滑らかな四次元空間である。
$\widetilde q:\mathcal H\to\mathcal L$を
$(p,a)\mapsto(p,r(a))$とすると
\[
 \widetilde D_{F,c}
 =\mathcal H\times_{\mathcal L}\widetilde Y_{F,c}.
\]
この写像は相対次元五の二次ファイブレーションで、
臨界部分およびその像はいずれも零切断$X$である。
解析局所またはétale局所で写像は
\[
 (u_1,u_2,u_3,z_1,\ldots,z_6)\longmapsto
 (u_1,u_2,u_3,z_1z_2+z_3z_4+z_5z_6)
\]
と同型である。
\end{proposition}
\begin{proof}
$x=tp$を$F(x)+cH_F(x)$に代入して$t^3$を除くと
$\widetilde Y_{F,c}$の方程式が得られる。
零切断では$dF\ne0$であり滑らか。
式に$t=r(a)$を代入するとファイバー積の等式が得られる。
各rank六の束ファイバー上で$r$は非退化二次形式である。
$r(a)=t$で定義されたファイバーの垂直微分が零となるのは
$a=0$のときだけで、同時に$t=0$を要求する。
滑らかな因子$X$に沿って$t$を法座標、
残りを局所座標$u_1,u_2,u_3$に取る。
非退化二次形式の局所標準化から表示した模型を得る。
\end{proof}
\begin{remark}[射影のファイバーのHessian因子]
$c\ne0$、$p\in X$に対し
$\widetilde D_{F,c}\to\PP^4$のファイバーは
$H_F(p)\ne0$なら$\{r^2=0\}\subset\CC^6$、
$H_F(p)=0$なら$\CC^6$である。
この射影の選択はPfaffian graphの模型に依存し、
元の対象モジュライ上の選択不要な不変量とはまだ主張しない。
\end{remark}

\begin{remark}[検証の状態]
ここでの数学的証明は、旧論文IIIの第23--24節に対応する
記述を現行TeXへ再構成したもので、今回の独立監査を意味しない。
一般の実際のBrill--Noether境界との接触同値、
とりわけ$F(q)+H_F(q)/96$への普遍的・解析的な正規化は
別の結果であり、本節では入力としても仮定なしに使用しない。
\end{remark}
